\chapter*{Résumé}
\addcontentsline{toc}{chapter}{Résumé}

Ce rapport présente l'analyse architecturale, logicielle et opérationnelle complète de la plateforme \textbf{Perceptra}, un système de recherche d'images par le contenu visuel (\textit{Content-Based Image Retrieval} -- CBIR) développé par le \textbf{Groupe~13} au sein de la Faculté des Sciences et Techniques de Tanger (FSTT) pour l'année universitaire 2026--2027.

Pour surmonter le \textit{fossé sémantique} (\textit{semantic gap}) sans dépendre d'annotations textuelles manuelles, Perceptra combine cinq descripteurs visuels déterministes~: histogrammes de couleur HSV (768-D), moments statistiques de Stricker \& Orengo (9-D), couleurs dominantes par partitionnement \textit{MiniBatchKMeans}, matrices de co-occurrence de niveaux de gris (GLCM de Haralick) et motifs binaires locaux (LBP 256-D).

L'apport algorithmique majeur du Groupe~13 réside dans la conception du moteur d'estimation de pertinence des descripteurs (\textbf{\textit{Feature Relevance Estimation} -- FRE}). En analysant en direct les rétroactions de l'utilisateur (images pertinentes et non-pertinentes) via un ratio de Fisher, une modélisation de dispersion intra-classe inverse et un lissage bayésien avec mémoire d'inertie, le système ajuste dynamiquement la pondération des caractéristiques tout en translatant le centroïde de la requête.

Sur le plan architectural, Perceptra associe un \textbf{Backend RESTful} en Python 3.13 sous Flask 3 structuré en couche de services découplés avec OpenCV et scikit-learn, un \textbf{Frontend réactif} en SvelteKit 3 et Svelte 5 Runes avec Shadcn-Svelte, Tailwind CSS 4 et Paraglide-JS, ainsi qu'une persistance hybride (fichiers vérifiés par SHA-256 et base relationnelle SQLAlchemy). Sur le jeu de référence \textbf{GHIM-10K} (10\,000 images), la précision au rang 10 ($P@10$) progresse de \textbf{58,2\,\%} à \textbf{85,1\,\%} après deux itérations de feedback FRE (+26,9 points de gain).

\vspace{0.2cm}
\noindent\textbf{Mots-clés~:} CBIR, Feature Relevance Estimation (FRE), Relevance Feedback, GLCM, LBP, Svelte 5 Runes, Flask, GHIM-10K, Architecture logicielle.

\vspace{0.3cm}
\begin{infobox}{Apport principal du Groupe 13}
Formalisation mathématique et réalisation logicielle d'un moteur d'estimation de pertinence des descripteurs (\textit{Feature Relevance Estimation}), transformant la recherche statique en un dialogue adaptatif convergent et élevant la précision de plus de 26 points sur GHIM-10K.
\end{infobox}

\newpage
\chapter*{Abstract}
\addcontentsline{toc}{chapter}{Abstract}

This report delivers a thorough architectural, algorithmic, and engineering analysis of \textbf{Perceptra}, an advanced Content-Based Image Retrieval (CBIR) platform integrating user-centric interactive learning through Relevance Feedback. Developed as the capstone project for Group~13 at the Faculty of Sciences and Techniques of Tangier (FSTT), Perceptra directly bridges the \textit{semantic gap} between low-level pixel data and human subjective visual concepts.

Rather than relying on manual textual metadata, Perceptra extracts a multi-faceted feature representation combining 768-D L1-normalized HSV color histograms, 9-D HSV statistical color moments, dominant color centroids via MiniBatchKMeans clustering, multi-directional Gray-Level Co-occurrence Matrices (GLCM), and 256-bin Local Binary Patterns (LBP).

Group~13's specialized contribution centers on the design and execution of the \textbf{Feature Relevance Estimation (FRE)} engine. By continuously assessing user feedback (positive and negative visual samples) via a Fisher-inspired discriminant ratio, intra-class variance modeling, and Bayesian smoothing, the retrieval system adaptively refines feature weights and shifts the query centroid towards the user's intended target subspace.

Architecturally, the application adopts a modern decoupled paradigm: a Python 3.13 Flask backend designed under the Service Layer pattern with OpenCV, scikit-image, and scikit-learn; a high-performance frontend powered by SvelteKit 3 and Svelte 5 Runes alongside Tailwind CSS 4 and Shadcn-Svelte; and a hybrid storage layer backed by SHA-256 content verification and SQLAlchemy relational persistence. Rigorous quantitative benchmarks on the 10,000-image GHIM-10K collection demonstrate an increase in Precision@10 from 58.2\% to 85.1\% after two feedback iterations.

\vspace{0.4cm}
\noindent\textbf{Keywords:} Content-Based Image Retrieval (CBIR), Feature Relevance Estimation (FRE), Interactive Relevance Feedback, Computer Vision, Texture Analysis, Svelte 5, Flask, GHIM-10K Dataset.
